\documentclass[../../../../SoftwareRequirementsSpecification.tex]{subfiles}
\begin{document}
% Richiede le macro definite in src/macros/useCaseMacros.tex
% (incluse nel preambolo di SoftwareRequirementsSpecification.tex)

\ucstart
\begin{tabularx}{\textwidth}{|l|X|}
  \hline
  \ucid{PT-04} \\ \hline
  \ucname{Configurazione Sessioni di Allenamento in una
    Scheda} \\ \hline
  \ucactor{PT} \\ \hline
  \ucdescription{Il PT inserisce o modifica gli
    allenamenti presenti in una scheda.} \\ \hline
  \ucprecond{Il PT deve essere loggato e deve
    essere presente almeno una scheda da modificare, che non abbia
    assegnazioni attive.} \\ \hline
  \ucpostcond{Vengono configurati gli allenamenti di
    una scheda.} \\ \hline
  \ucmainflow{
    \item Il PT atterra sulla pagina di configurazione delle
      sessioni di allenamento.
    \item Il sistema mostra la lista di sessioni di allenamento presenti.
    \item Il PT preme su una delle sessioni di allenamento.
    \item Il sistema mostra il dettaglio della sessione di allenamento corrente.
    \item Il PT preme il tasto "Aggiungi Esercizio".
    \item Il sistema mostra una lista di esercizi.
    \item Il PT sceglie un esercizio dalla lista.
    \item Il sistema mostra il form di inserimento dettagli esercizio.
    \item Il PT configura l'esercizio inserendo:
      \begin{itemize}
        \item Numero di serie/ripetizioni.
        \item Tempo di recupero tra una serie e l'altra (in secondi).
        \item Variante di esecuzione se presente (superserie, dropset,
          rest-pause etc...).
        \item Tipo di progressione se presente (di carico, di volume,
          doppia progressione o di densità).
        \item Altre note esecutive.
      \end{itemize}
    \item Il PT preme il tasto "Aggiungi Esercizio".
    \item Il sistema aggiunge l'esercizio alla sessione.
    \item Il flusso riprende dal punto 4.
  } \\ \hline
  \ucaltflow{
    \textbf{5a. Modifica esercizio} \newline
    - Invece di aggiungere un nuovo esercizio il PT preme su
    un'esercizio già presente e modifica quello. \newline
    - Il flusso riprende dal punto 4. \newline
    \textbf{5b. Modifica sessione scelta} \newline
    - Il PT decide di modificare un'altra sessione di
    allenamento invece di quella corrente.\newline
    - Il flusso riprende dal punto 1.\newline
    \textbf{9c. Progressione non compatibile con l'esercizio} \newline
    - Non tutte le progressioni sono applicabili a tutti gli esercizi
    (una progressione di carico non ha senso su un esercizio a corpo
    libero, una di volume su un esercizio a cedimento). \newline
    - Il sistema propone quindi solo le progressioni compatibili con
    l'esercizio scelto al passo 7; se nessuna lo è, il campo non
    viene proposto. \newline
    - Il flusso riprende dal punto 9.
  } \\ \hline
  \ucvariations{
    \textbf{9a. Selezione variante di esecuzione} \newline
    - Se viene selezionata una variante potrebbero essere richiesti dati
    aggiuntivi (es. altri esercizi da effettuare in superserie, tempo di recupero
    tra una miniserie e l'altra nel rest-pause etc..).\newline
    \textbf{9b. Selezione tipo di progressione}\newline
    - Se viene selezionata una progressione potrebbero essere richiesti
    dati aggiuntivi (es. modifica settimanale del peso nella
      progressione di carico, delle serie/ripetizioni nella
      progressione di volume, di entrambi nella doppia progressione, o
    del tempo di recupero in quella di densità.)
  } \\ \hline
\end{tabularx}
\end{document}
